% !TEX root = ../monograph.tex
\subsection{Obtaining the three comparison bounds from one lattice family}
\label{ch14:lattice-comparison}

Here are the scale choices behind \eqref{ch14:grid-input}. Fix the
upper bound $n$ on the number of insertions, maximum scaling dimension $D$ and a compact set of configurations
$K$ with positive point separation. Enlarge the finite label list to
include its signed coordinate images and conjugates. All constants and
all three exponents below are chosen for this one finite request.
Choose positive numbers $\alpha_0,\alpha_1,\alpha_2$ and $c<0$ with
\begin{align}
 \alpha_1+6&<\alpha_0,&\alpha_2+5&<\alpha_0+c,\notag\\
 nD(\alpha_0-\alpha_1)+c&<0,&
 nD(\alpha_0-\alpha_2)-\alpha_0-c&<0.
 \label{ch14:comparison-exponents}
\end{align}
For example first choose $\alpha_0-\alpha_1=7$, next $-c>7nD$,
then $\alpha_0-\alpha_2>5-c$, and finally increase all three
$\alpha_i$ together until the last inequality holds and all are
positive. Increasing their common part does not alter the preceding
inequalities.

At a comparison scale $k$ whose comparison window lies in a single stage put\mgubi{Issue 10: say for which scales this holds ($k=k_j$ suffices for the limit taken below) and cite Chapter~\ref{ch10:trajectory}.}
$H_k=\exp((\log\log\widehat\beta_k)^2)$ and set
\[
 r_k=a_kH_k^{\alpha_0},\qquad s_k=H_k^{-\alpha_0},\qquad
 \delta_k=H_k^c.
\]
Thus $r_ks_k=a_k$: grid rounding occurs on the construction
lattice. Use one common localization map $T$ from
Section~\ref{ch11:common-extraction}. If $v$ is a fixed real polynomial
representative of a free quotient field, its lift $u_{k,N,p}$ solves
\[
                 S_{g/\theta}T u_{k,N,p}=W_{H_k^{\alpha_0}}v.
\]
Here $S_t$ multiplies a monomial with $f$ curvature factors by $t^f$,
whereas $W_\rho$ multiplies one of scaling dimension $j$ by $\rho^j$;
$g=\beta_{N,p,k}^{-1/2}$ and $\theta$ is the local normalization
used for $T$. These two dilations have different purposes.
The inverse bound in \eqref{ch11:simultaneous-detector-lifts} gives
bounded lifts at each fixed $k$, uniformly in the cutoffs and periods.
Their identities for reality, the identity field and signed symmetries are exact. No separate
choice is made for different spectator words.

For clarity, the use for finitely many insertions of the estimate for the corrected coefficients is
as follows. Write $I_i$ for the lattice insertion in argument $i$. For a subset
$J$ of the insertion indices, let $G_J$ be the free correlation of the
corresponding fields at the rescaled positions. Define recursively
\[
 R_\varnothing=1,\qquad
 R_I=\prod_{i\in I}I_i-
          \sum_{\varnothing\ne J\subseteq I}G_JR_{I\setminus J}.
\]
These are the Wick products of the insertions with respect to the free field (the moments $G_J$ are those of the free field, not of the lattice measure).
Equivalently, introduce commuting variables $s_i$ with $s_i^2=0$ and write
$s_J=\prod_{i\in J}s_i$. Then $R_I$ is the coefficient of $s_I$ in
$\prod_i(1+s_i I_i)/(1+\sum_{J\ne\varnothing}G_Js_J)$. The denominator has
constant coefficient one and hence an exact inverse in the finite
squarefree algebra. For the limiting Gaussian variables every
nonempty coefficient has expectation zero. The extraction of coefficients of source monomials,
covariance comparison and absorption of the larger source weight of
Proposition~\ref{ch11:n-mark-comparison} apply to this \emph{entire}
finite algebra, including its joint incoming terms.
With $q=\log H_k/(16n)$ and source weight $e^{2q}$, the output
normalization is at most $e^q$ per reflected cell, and extraction gives
\begin{equation}
 \|R_I\|_{\mathrm{cell}}
       \le e^{q-2q|I|}\le H_k^{-1/(16n)}
       \quad(\varnothing\ne I\subseteq\{1,\ldots,n\}).
 \label{ch14:subset-root}
\end{equation}
The lower bound on the reference coupling at the stage is chosen after this fixed number of insertions and the constants of its
insertions. The estimate is on the original positive partition
function and includes every reflected copy. Uniformity on $K$ uses
positive point separation in the covariance estimate; it does not
require a positive lower bound on each coordinate difference.
Unwinding the finite recursion and applying the chessboard inequality
therefore gives the grid error $C_KH_k^{-1/(16n)}$ at generic grid
points. Distinct coordinates there differ by at least $a_k$ in continuum
units, so the assignments to lattice sites require no rounding.

Apply the same coefficient map and the same finite comparison request
also at exponents $\alpha_1$ and $\alpha_2$. At its own normalization
a single insertion has chessboard seminorm at most $1+|G_{\{i\}}|$. Changing its
scaling dilation to $\alpha_0$ loses at most a factor
$CH_k^{(\alpha_0-\alpha_j)D}$. Set
\[
 A_k=CH_k^{(\alpha_0-\alpha_1)D},\qquad
 V_k=CH_k^{(\alpha_0-\alpha_2)D}.
\]
The comparison cells have sides at most $a_kH_k^{\alpha_j+5}$.
Consequently the first cells are $o(r_k)$, and the second are
$o(r_k\delta_k)$, by \eqref{ch14:comparison-exponents}.
The chessboard inequality gives $|S_k|\le A_k^n$ on $K$.
The modulus of continuity in the insertion points of Section~\ref{ch11:positional}, applied
with margin $r_k\delta_k$ in continuum units, gives
$|\nabla S_k|\le4nV_k^n/\delta_k$ outside the coincident-coordinate strips.
Include a strip for each configuration taking every coordinate from the point or
from its rounding (finitely many); the total volume of the strips is still $O(\delta_k)$.
The compact configuration sets in continuum units lie inside the smallest period for large $k$,
so the periodic and ordinary separations agree.

Take an ordered limit point of the lifts at each fixed $k$. The
uniform bounds pass to the already chosen continuum hierarchy.
Their exponents give exactly \eqref{ch14:grid-cost}, since
\[
 A_k^n\delta_k=CH_k^{nD(\alpha_0-\alpha_1)+c},\qquad
 V_k^ns_k/\delta_k=CH_k^{nD(\alpha_0-\alpha_2)-\alpha_0-c}.
\]
Both tend to zero, and $s_k/\delta_k\to0$ follows from the second
inequality. A diagonal sequence over the finite requests gives a
single comparison sequence. The bounded coefficient limits do not
alter the chosen correlation hierarchy or its order of limits.
